\section{讨论}
\label{sec:discussion}
\subsection{结果支持的结论}
主要证据支持在本研究取回总体中使用操作需求作为空中感知的条件。
受控差异是空间加权：相同采集模型与候选可以将测量导向支撑未来机械臂站位的区域，
而不仅导向通用缺失度。无已确认候选失败的减少与这一解释一致，
而完整物理终点避免将候选发现误当成成功取回。

相关性场不应理解为成功概率。它组合经验证样本、关节余量、占地轮廓与当前操作证据，
但没有建模全部 IK 分支，也没有联合优化导航、相机可见性和抓取执行。
单步增益保持当前支撑固定，而后续证据可能阻断或揭示其他站位。
包括仅 Generic 成功场景在内的剩余失败表明，相关性引导不必在每个场景中占优。
本文不确立单调任务进展或全局最优感知保证。

\subsection{任务、感知与仿真边界}
\label{sec:limitations}
已评价总体包含一个已知砖块、静态水平地面、当前机器人几何，以及零/一/两墙布局的有限扰动。
初始感知从已知作业区域开始，不是对任意物体的无限范围空中搜索。
RGB-D 假设、地面任务可达性图、平面参考桥接与夹爪几何均与任务相关。
不平地形、动态遮挡、不同物体几何及未见布局均超出当前证据范围。

有限扫描机会在名义悬停位姿下对确定扫描计划的起始相位取平均，
并未校准真实回波概率、反射率、位姿误差或点云内部运动畸变。
假设高度遮挡体与未知空间乐观可见性可能误判下一窗口的实际观测。
环境证据仍包含粗粒度保守阻断，有限候选验证也可能漏掉有用站位。
预算内缺少真实地面支撑仍是合理停止原因；预测机会不能替代测量。

物理执行比运动学可达计数更接近任务终点，但仍不是实机验证。
当前平台的无人机公共 map 对齐使用私有仿真机体状态后端，地面全局对齐由出生位姿派生，
抓取确认使用仿真双侧接触与关节反馈。
策略输入是运行时传感器和公共变换，但完整定位/接触栈并非真实世界意义上的纯传感器链路。
独立检查器还使用仿真物体高度验证提升。
短时保持由任务的抓取确认保持支持，本研究不提供覆盖长时间保持的独立连续物体高度轨迹。
动力学、摩擦、标定和反馈差异限制了结果迁移。

\subsection{实验解释与未来迁移}
96 个配对场景支持预定主对比，而非普遍泛化结论。
难度是抽样混合，定性案例由结果选择，辅助子集较小。
条件效率只包含共同成功任务，距离缺测使距离节省主张不成立。
因此，配对成功率结果应保持为核心经验结论，不应扩展为未经测量的能耗、速度或实机稳健性优势。

未来迁移应先建立真实共享 map/时间约定、相机与雷达外参、手眼标定，
再将公共飞行、底盘和轨迹接口绑定到实测硬件反馈。
真实夹持与物体保持需要独立表征，仅控制器完成不足以证明成功。
静态传感/反馈数据和离线几何检查适合作为有界运动测试前的第一步。
本文不包含真实适配器或硬件实验；当前仿真版本和结果将作为后续迁移研究的固定参照。
